%--------------------------------------------------------------------------------
%
%
%--------------------------------------------------------------------------------
\pagebreak
\section*{RFI — Return from Interrupt}
\addcontentsline{toc}{section}{RFI — Return from Interrupt}

\begin{iPageDescEntry}{Syntax}
    \texttt{RFI [ RegR ]}
\end{iPageDescEntry}

\begin{iPageDescEntry}{Format}
    The \texttt{RFI} instruction uses the branch instruction format. \\

    \begin{flushleft}
    \drawInstrFormat
        {0,1,3,5,6.5, 16}
        {\OpCodeGrpSYS, \OpCodeRFI, RegR, 0, 0 }
    \end{flushleft}
\end{iPageDescEntry}

\begin{iPageDescEntry}{Description}
    The \texttt{RFI} instruction restores processor state after an interrupt
    or exception handler has completed.  

    Execution resumes at the program counter saved during the interrupt,
    and the processor state (including privilege level, interrupt enable flags,
    and other PSR fields) is restored from the saved copy.
\end{iPageDescEntry}

\begin{iPageDescEntry}{Operation}
    \texttt{PSW ← IPSW;} 
\end{iPageDescEntry}

\begin{iPageDescEntry}{Exceptions}
    None.
\end{iPageDescEntry}

\begin{iPageDescEntry}{Notes}
    - Only valid in supervisor mode. \\
    - \texttt{RegR} may optionally hold the return address if software requires. \\
    - Typically paired with \texttt{INT} / \texttt{TRAP} entry sequences.
\end{iPageDescEntry}



